\documentclass[a4paper,12pt,twoside]{memoir}

% Castellano
\usepackage[spanish,es-tabla]{babel}
\selectlanguage{spanish}
\usepackage[utf8]{inputenc}
\usepackage[T1]{fontenc}
\usepackage{lmodern} % Scalable font
\usepackage{microtype}
\usepackage{placeins}

\RequirePackage{booktabs}
\RequirePackage[table]{xcolor}
\RequirePackage{xtab}
\RequirePackage{multirow}

% Links
\PassOptionsToPackage{hyphens}{url}\usepackage[colorlinks]{hyperref}
\hypersetup{
	allcolors = {red}
}

% Ecuaciones
\usepackage{amsmath}

% Rutas de fichero / paquete
\newcommand{\ruta}[1]{{\sffamily #1}}

% Párrafos
\nonzeroparskip

% Huérfanas y viudas
\widowpenalty100000
\clubpenalty100000

% Imagenes
\usepackage{graphicx}
\newcommand{\imagen}[2]{
	\begin{figure}[!h]
		\centering
		\includegraphics[width=0.9\textwidth]{#1}
		\caption{#2}\label{fig:#1}
	\end{figure}
	\FloatBarrier
}

\newcommand{\imagenflotante}[2]{
	\begin{figure}%[!h]
		\centering
		\includegraphics[width=0.9\textwidth]{#1}
		\caption{#2}\label{fig:#1}
	\end{figure}
}



% El comando \figura nos permite insertar figuras comodamente, y utilizando
% siempre el mismo formato. Los parametros son:
% 1 -> Porcentaje del ancho de página que ocupará la figura (de 0 a 1)
% 2 --> Fichero de la imagen
% 3 --> Texto a pie de imagen
% 4 --> Etiqueta (label) para referencias
% 5 --> Opciones que queramos pasarle al \includegraphics
% 6 --> Opciones de posicionamiento a pasarle a \begin{figure}
\newcommand{\figuraConPosicion}[6]{%
  \setlength{\anchoFloat}{#1\textwidth}%
  \addtolength{\anchoFloat}{-4\fboxsep}%
  \setlength{\anchoFigura}{\anchoFloat}%
  \begin{figure}[#6]
    \begin{center}%
      \Ovalbox{%
        \begin{minipage}{\anchoFloat}%
          \begin{center}%
            \includegraphics[width=\anchoFigura,#5]{#2}%
            \caption{#3}%
            \label{#4}%
          \end{center}%
        \end{minipage}
      }%
    \end{center}%
  \end{figure}%
}

%
% Comando para incluir imágenes en formato apaisado (sin marco).
\newcommand{\figuraApaisadaSinMarco}[5]{%
  \begin{figure}%
    \begin{center}%
    \includegraphics[angle=90,height=#1\textheight,#5]{#2}%
    \caption{#3}%
    \label{#4}%
    \end{center}%
  \end{figure}%
}
% Para las tablas
\newcommand{\otoprule}{\midrule [\heavyrulewidth]}
%
% Nuevo comando para tablas pequeñas (menos de una página).
\newcommand{\tablaSmall}[5]{%
 \begin{table}
  \begin{center}
   \rowcolors {2}{gray!35}{}
   \begin{tabular}{#2}
    \toprule
    #4
    \otoprule
    #5
    \bottomrule
   \end{tabular}
   \caption{#1}
   \label{tabla:#3}
  \end{center}
 \end{table}
}

%
% Nuevo comando para tablas pequeñas (menos de una página).
\newcommand{\tablaSmallSinColores}[5]{%
 \begin{table}[H]
  \begin{center}
   \begin{tabular}{#2}
    \toprule
    #4
    \otoprule
    #5
    \bottomrule
   \end{tabular}
   \caption{#1}
   \label{tabla:#3}
  \end{center}
 \end{table}
}

\newcommand{\tablaApaisadaSmall}[5]{%
\begin{landscape}
  \begin{table}
   \begin{center}
    \rowcolors {2}{gray!35}{}
    \begin{tabular}{#2}
     \toprule
     #4
     \otoprule
     #5
     \bottomrule
    \end{tabular}
    \caption{#1}
    \label{tabla:#3}
   \end{center}
  \end{table}
\end{landscape}
}

%
% Nuevo comando para tablas grandes con cabecera y filas alternas coloreadas en gris.
\newcommand{\tabla}[6]{%
  \begin{center}
    \tablefirsthead{
      \toprule
      #5
      \otoprule
    }
    \tablehead{
      \multicolumn{#3}{l}{\small\sl continúa desde la página anterior}\\
      \toprule
      #5
      \otoprule
    }
    \tabletail{
      \hline
      \multicolumn{#3}{r}{\small\sl continúa en la página siguiente}\\
    }
    \tablelasttail{
      \hline
    }
    \bottomcaption{#1}
    \rowcolors {2}{gray!35}{}
    \begin{xtabular}{#2}
      #6
      \bottomrule
    \end{xtabular}
    \label{tabla:#4}
  \end{center}
}

%
% Nuevo comando para tablas grandes con cabecera.
\newcommand{\tablaSinColores}[6]{%
  \begin{center}
    \tablefirsthead{
      \toprule
      #5
      \otoprule
    }
    \tablehead{
      \multicolumn{#3}{l}{\small\sl continúa desde la página anterior}\\
      \toprule
      #5
      \otoprule
    }
    \tabletail{
      \hline
      \multicolumn{#3}{r}{\small\sl continúa en la página siguiente}\\
    }
    \tablelasttail{
      \hline
    }
    \bottomcaption{#1}
    \begin{xtabular}{#2}
      #6
      \bottomrule
    \end{xtabular}
    \label{tabla:#4}
  \end{center}
}

%
% Nuevo comando para tablas grandes sin cabecera.
\newcommand{\tablaSinCabecera}[5]{%
  \begin{center}
    \tablefirsthead{
      \toprule
    }
    \tablehead{
      \multicolumn{#3}{l}{\small\sl continúa desde la página anterior}\\
      \hline
    }
    \tabletail{
      \hline
      \multicolumn{#3}{r}{\small\sl continúa en la página siguiente}\\
    }
    \tablelasttail{
      \hline
    }
    \bottomcaption{#1}
  \begin{xtabular}{#2}
    #5
   \bottomrule
  \end{xtabular}
  \label{tabla:#4}
  \end{center}
}



\definecolor{cgoLight}{HTML}{EEEEEE}
\definecolor{cgoExtralight}{HTML}{FFFFFF}

%
% Nuevo comando para tablas grandes sin cabecera.
\newcommand{\tablaSinCabeceraConBandas}[5]{%
  \begin{center}
    \tablefirsthead{
      \toprule
    }
    \tablehead{
      \multicolumn{#3}{l}{\small\sl continúa desde la página anterior}\\
      \hline
    }
    \tabletail{
      \hline
      \multicolumn{#3}{r}{\small\sl continúa en la página siguiente}\\
    }
    \tablelasttail{
      \hline
    }
    \bottomcaption{#1}
    \rowcolors[]{1}{cgoExtralight}{cgoLight}

  \begin{xtabular}{#2}
    #5
   \bottomrule
  \end{xtabular}
  \label{tabla:#4}
  \end{center}
}


















\graphicspath{ {./img/} }

% Capítulos
\chapterstyle{bianchi}
\newcommand{\capitulo}[2]{
	\setcounter{chapter}{#1}
	\setcounter{section}{0}
	\chapter*{#2}
	\addcontentsline{toc}{chapter}{#2}
	\markboth{#2}{#2}
}

% Apéndices
\renewcommand{\appendixname}{Apéndice}
\renewcommand*\cftappendixname{\appendixname}

\newcommand{\apendice}[1]{
	%\renewcommand{\thechapter}{A}
	\chapter{#1}
}

\renewcommand*\cftappendixname{\appendixname\ }

% Formato de portada
\makeatletter
\usepackage{xcolor}
\newcommand{\tutor}[1]{\def\@tutor{#1}}
\newcommand{\course}[1]{\def\@course{#1}}
\definecolor{cpardoBox}{HTML}{E6E6FF}
\def\maketitle{
  \null
  \thispagestyle{empty}
  % Cabecera ----------------
\noindent\includegraphics[width=\textwidth]{cabecera}\vspace{1cm}%
  \vfill
  % Título proyecto y escudo informática ----------------
  \colorbox{cpardoBox}{%
    \begin{minipage}{.8\textwidth}
      \vspace{.5cm}\Large
      \begin{center}
      \textbf{TFG del Grado en Ingeniería Informática}\vspace{.6cm}\\
      \textbf{\LARGE\@title{}}
      \end{center}
      \vspace{.2cm}
    \end{minipage}

  }%
  \hfill\begin{minipage}{.20\textwidth}
    \includegraphics[width=\textwidth]{escudoInfor}
  \end{minipage}
  \vfill
  % Datos de alumno, curso y tutores ------------------
  \begin{center}%
  {%
    \noindent\LARGE
    Presentado por \@author{}\\ 
    en Universidad de Burgos --- \@date{}\\
    Tutor: \@tutor{}\\
  }%
  \end{center}%
  \null
  \cleardoublepage
  }
\makeatother

\newcommand{\nombre}{Nombre del alumno} %%% cambio de comando

% Datos de portada
\title{título del TFG}
\author{\nombre}
\tutor{nombre tutor}
\date{\today}

\begin{document}

\maketitle


\newpage\null\thispagestyle{empty}\newpage


%%%%%%%%%%%%%%%%%%%%%%%%%%%%%%%%%%%%%%%%%%%%%%%%%%%%%%%%%%%%%%%%%%%%%%%%%%%%%%%%%%%%%%%%
\thispagestyle{empty}


\noindent\includegraphics[width=\textwidth]{cabecera}\vspace{1cm}

\noindent D. nombre tutor, profesor del departamento de nombre departamento, área de nombre área.

\noindent Expone:

\noindent Que el alumno D. \nombre, con DNI dni, ha realizado el Trabajo final de Grado en Ingeniería Informática titulado título de TFG. 

\noindent Y que dicho trabajo ha sido realizado por el alumno bajo la dirección del que suscribe, en virtud de lo cual se autoriza su presentación y defensa.

\begin{center} %\large
En Burgos, {\large \today}
\end{center}

\vfill\vfill\vfill

% Author and supervisor
\begin{minipage}{0.45\textwidth}
\begin{flushleft} %\large
Vº. Bº. del Tutor:\\[2cm]
D. nombre tutor
\end{flushleft}
\end{minipage}
\hfill
\begin{minipage}{0.45\textwidth}
\begin{flushleft} %\large
Vº. Bº. del co-tutor:\\[2cm]
D. nombre co-tutor
\end{flushleft}
\end{minipage}
\hfill

\vfill

% para casos con solo un tutor comentar lo anterior
% y descomentar lo siguiente
%Vº. Bº. del Tutor:\\[2cm]
%D. nombre tutor


\newpage\null\thispagestyle{empty}\newpage




\frontmatter

% Abstract en castellano
\renewcommand*\abstractname{Resumen}
\begin{abstract}
En este primer apartado se hace una \textbf{breve} presentación del tema que se aborda en el proyecto.
\end{abstract}

\renewcommand*\abstractname{Descriptores}
\begin{abstract}
Palabras separadas por comas que identifiquen el contenido del proyecto Ej: servidor web, buscador de vuelos, android \ldots
\end{abstract}

\clearpage

% Abstract en inglés
\renewcommand*\abstractname{Abstract}
\begin{abstract}
A \textbf{brief} presentation of the topic addressed in the project.
\end{abstract}

\renewcommand*\abstractname{Keywords}
\begin{abstract}
keywords separated by commas.
\end{abstract}

\clearpage

% Indices
\tableofcontents

\clearpage

\listoffigures

\clearpage

\listoftables
\clearpage

\mainmatter
\capitulo{1}{Introducción}

La gestión de las reservas de las aulas siempre ha sido un problema en el que se ven envueltos los profesores que quieren reservar y los encargados de cada aula en cuestión.
Esta aplicación web da solución a este problema, ya que permite la visualización de eventos y aulas a todos los usuarios y la posibilidad de realizar reservas de forma sencilla para usuarios que tengan cuenta de \textit{Office 365} o Microsoft (gestionado por el login de Microsoft que utiliza \textit{OAuth2} como estándar de autorización) y los permisos necesarios de las aulas sin que existan solapamientos entre dichas reservas. En este proyecto los calendarios, los eventos y las horas se gestionarán desde la API REST de Microsoft, en concreto utilizando el punto de enlace de \textit{Microsoft Graph}.
\capitulo{2}{Objetivos del proyecto}

Este apartado explica de forma precisa y concisa cuales son los objetivos que se persiguen con la realización del proyecto. Se puede distinguir entre los objetivos marcados por los requisitos del software a construir y los objetivos de carácter técnico que plantea a la hora de llevar a la práctica el proyecto.

\capitulo{3}{Conceptos teóricos}

\subsection{Frameworks de desarrollo}
Un framework es un entorno o marco de trabajo estandarizado, con una estructura conceptual y tecnológica que puede servir de base para la organización y desarrollo de software \cite{wiki:framework}.\newline
En este proyecto he utilizado como framework Flask, que se explicará en las herramientas usadas.

\subsection{API REST}
"Una API REST es cualquier interfaz entre sistemas que utilice directamente una llamada de HTTP para obtener datos o indicar la ejecución de operaciones sobre los datos, en cualquier formato (XML, JSON, etc)". Esto permite que una aplicación puede interactuar con un recurso proporcionado en REST conociendo el identificador del recurso y la acción requerida para obtener la finalidad deseada \cite{wiki:REST}.\newline
En nuestro caso se ha utilizado el endpoint de Microsoft Graph para acceder a la API REST de Outlook, para ello hubo que realizar lo siguiente:\newline
Primeramente se utiliza la API de Microsoft para  para realizar el login en la aplicación mediante una cuenta de Outlook. Utiliza OAuth 2.0, que necesita el token recibido al registrar la aplicación en Azure Active Directory y nos devuelve un nuevo token. Una vez lo recibimos podemos hacer uso del endpoint de Microsoft Graph mediante este token de OAuth recibido.
\newline
Una vez se dispone de este acceso, se realizan llamadas mediante el punto de conexión Microsoft Graph: punto de conexión común para acceder a la API REST que permite el acceso y la modificación a todas las características de una cuenta Microsoft O365. Dentro de este punto de conexión se encuentran las llamadas al resto de las APIs de Microsoft para lo que necesitemos (calendarios, usuarios, grupos de calendarios, correo).


\subsection{APIs de validación}
Permite implementar el inicio de sesión que utiliza una aplicación en tu propio proyecto, de forma que si existe esa cuenta permite el acceso a la dirección que se haya especificado. Para que este proceso funcione, Microsoft utiliza OAuth v2.0, que sigue un flujo de autorización y autenticación para acceder a la aplicación y redireccionar al usuario a la página indicada en Azure Active Directory. Este proceso devuelve un token de acceso a la aplicación que sirve para obtener acceso a los recursos protegidos mediante el punto de conexión de la API utilizada y otro token para refrescar el código de acceso anterior en caso de que caduque si el inicio de sesión es correcto. Un token es un código de letras y números único que se necesita para realizar el resto de llamadas a la API en la que nos hemos logeado para obtener información sobre la sesión que hemos iniciado.

\subsection{APIs de calendarios}
La API de un calendario nos permite guardar eventos dentro de los calendarios que nos proporciona, borrar, modificar o compartir estos eventos y calendarios. Mediante llamadas a esta API se pueden acceder a los calendarios, recuperar los eventos y realizar las distintas modificaciones sobre ellos. Esta API que hemos usado pertenece a la API REST de Outlook y accedemos a ella mediante el punto de acceso de Microsoft Graph como se ha explicado anteriormente.

\subsection{APIs de conectividad de BD}
Es una capa intermedia de acceso a las bases de datos que hace posible las acciones \textit{CRUD} a cualquier dato desde la aplicación que desarrollamos \cite{wiki:ODBC}. En en el presente trabajo utilizamos \textit{pyodbc}: Es un módulo de Python de código abierto que simplifica el acceso a las bases de datos ODBC. Usado para realizar la conexión entre la base de datos y la aplicación web desarrollada, alojadas ambas en Azure.


\capitulo{4}{Técnicas y herramientas}

\section{Técnicas}

\subsection{\textit{Scrum}}
La metodología \textit{Scrum} es una metodología de desarrollo ágil, en la que se aplican buenas prácticas y procesos con la finalidad de obtener la mejor calidad posible en el producto final buscado.\newline
Esta metodología está pensada para un equipo de trabajo, con el fin de organizarse y trabajar más eficientemente obteniendo los mejores resultados posibles. En mi caso trataré de simular una metodología \textit{Scrum} ya que el trabajo le realizaré yo solo, pero la intención es dividir dicho trabajo en \textit{sprints} que contendrán dentro las tareas a realizar.\newline
Como acabo de mencionar, esta metodología se basa en crear \textit{sprints}, dentro de los que se crearán las tareas a realizar (\textit{issues}). Para establecer estas tareas dentro de cada sprint, se realiza una planificación, en la que se deja claro los objetivos a conseguir y el tiempo para cada uno de ellos. 
Se realizan revisiones para controlar los avances de los sprint planificados y llevar un control sobre el proyecto.

\subsubsection{\textit{GitHub}}
Para el control de versiones de este proyecto he utilizado \textit{GitHub}, que es un repositorio en línea que emplea \textit{Git}. Así podemos guardar los cambios que vamos realizando en el proyecto y controlar que se cumplan las tareas en el tiempo indicado.

\section{Herramientas}
\subsection{\textit{Visual Studio Code}}
Editor y compilador de código fuente. Incluye y utilizamos el soporte que ofrece para la depuración y para el control integrado de \textit{Azure Web Services}.

\subsection{\textit{Flask}}
\textit{Flask} es un \textit{framework} escrito en \textit{Python} que permite crear aplicaciones web rápidamente y con un mínimo número de líneas de código.

\subsection{\textit{Jinja2}}
\textit{Jinja2} es un motor de plantillas web para aplicaciones desarrolladas en \textit{Python}. Permite que \textit{Flask} pueda hacer uso de los contenidos de las plantillas \textit{HTML}.

\subsection{Documentación}
Como herramienta para realizar la documentación se ha escogido \textit{LaTeX}, está diseñado para crear documentos con una alta calidad tipográfica. Como editor de \textit{LaTeX} he utilizado \texit{Overleaf}

\subsection{Outlook API (Microsoft Graph) }
Para realizar este proyecto se ha utilizado la \texit{API} de \texit{Outlook} para manejar los calendarios, que equivalen a las aulas que se proporcionan al usuario para reservar. Para tener acceso a las herramientas de esta API, usé el punto de conexión Microsoft Graph ya que ofrece más servicios y características siguiendo la recomendación de Microsoft (\url{https://docs.microsoft.com/es-ES/outlook/rest/compare-graph}).
En concreto para este proyecto he utilizado distintas funciones de esta \texit{API}:
\begin{itemize}
    \item \texbf{Microsoft graph} para realizar las llamadas a esta \texit{API} y recoger o escribir datos en los calendarios, y modificar y crear dichos calendarios.
    \item \texbf{Outlook Login} para el inicio de sesión mediante una cuenta de \texit{Outlook}.
\end{itemize}

\capitulo{5}{Aspectos relevantes del desarrollo del proyecto}

En este apartado mostraré los aspectos más interesantes a la hora de desarrollar el proyecto, con esto me refiero a problemas surgidos, los caminos tomados para avanzar, el desarrollo del diseño y de la funcionalidad.

\subsection{Comienzo del proyecto}
Este proyecto me pareció útil e interesante, ya que se centra en el desarrollo web de una aplicación bastante realista. El desarrollo de una aplicación web no ha sido un tema que se haya visto mucho en la carrera por lo que me parecía interesante ya que en las prácticas trabajé un poco con ello y me gustó, por lo que me decidí a escoger este trabajo.   

\subsection{Preparación}
Como ya he comentado en la sección anterior, no tenía los conocimientos suficientes para el desarrollo de esta aplicación web, sobretodo tratando con lenguaje Python, exactamente con el \textit{framework Flask}. Sobre \textit{HTML} y \textit{JavaScript} ya estaba más relacionado.También tuve que tratar con la Microsoft Graph, desde dónde se gestionarían los eventos y calendarios.\newline
Para aprender sobre estos temas antes de comenzar seguí los siguientes tutoriales:\newline

\begin{itemize}
\item \href{https://blog.miguelgrinberg.com/post/the-flask-mega-tutorial-part-i-hello-world}{The Flask Mega Tutorial}\cite{flaskMegaTutorial}
\item \href{https://j2logo.com/tutorial-flask-espanol/}{J2logo}\cite{j2logo}
\end{itemize}
Cuanto más avanzaba el proyecto, más herramientas y distintas técnicas hacían falta, por lo que se tuvo que usar mucha documentación para comprender los funcionamientos de estas, las más consultadas han sido:
\begin{itemize}
\item \href{https://docs.microsoft.com/es-es/graph/api/overview?view=graph-rest-1.0}{Microsoft Graph}\cite{microsoftGraph}
\item \href{https://flask.palletsprojects.com/en/1.1.x/}{Documentación Flask}\cite{Flask}
\item \href{https://getbootstrap.com/docs/4.5/getting-started/introduction/}{Bootstrap}\cite{Bootstrap}
\end{itemize}

\subsection{Metodología del proyecto - Resumen}
Para comenzar el proyecto en Flask se necesitó tener instalado Visual Studio Code desde donde programaríamos nuestra web app, para ello hay que realizar unas instalaciones de requisitos que iban a ser necesarias para el desarrollo del proyecto desde el VSCode:
\begin{enumerate}
    \item Creación de un entorno virtual, para aislar las dependencias que requiere el proyecto de las ya existentes.
    \item Pylint como verificador del código que escribimos.
    \item Git Tortoise para facilitar las subidas a GitHub.
    \item Instalar extensiones de VSCode.
    \item Registrar la aplicación en Azure Active Directory.
\end{enumerate}
El primer punto a comentar es el uso de la API de Outlook, como he comentado accediendo a ella mediante Microsoft Graph, pero para poder acceder a los datos de Outlook y crear un inicio de sesión, primeramente tuve que registrarla en Azure Active Directory:
\begin{itemize}
    \item Acceder al centro de administración de Azure Active Directory.
    \item Acceder al apartado registro de aplicaciones y crear un nuevo registro.
    \item Una vez ahí dentro, se pedirán rellenar unos campos con el nombre de la aplicación, cuentas que accederán y lo más importante, la dirección URL de redirección al iniciar sesión mediante Outlook.
    \item Una vez se ha creado, recogemos el ID de la aplicación obtenido, que lo utilizaremos para obtener el secreto del cliente, que solamente se mostrará una vez. Para obtenerlo accedemos a certificados y secretos, se escribe un nombre y elige una duración y se obtiene dicho secreto.
\end{itemize}
Esto está especificado en el \href{https://docs.microsoft.com/es-es/graph/tutorials/python}{tutorial de Microsoft}\cite{pythonMicrosoftGraph}. \newline
Una vez registrada la aplicación tuve que seguir una serie de pasos para incluirlo en el código y hacer que funcionara:
\begin{itemize}
    \item Agregar al código la autenticación de Azure que hemos conseguido en los pasos anteriores (ver figura \ref{fig:AzureAutenticacion}) .
    \imagen{AzureAutenticacion}{Código recibido de Azure Active Directory}
    \item Agregar al código los permisos de Outlook que se van a proporcionar al usuario que acceda.
    \item Crear unas funciones con llamadas a la API REST para obtener un token de inicio de sesión y otro para refrescarle. Esto servirá para cada sesión de cada usuario diferente.
\end{itemize}
Una vez se creó este primer código que permitía acceder mediante el inicio de sesión de Outlook a la página que especificamos en Azure Active Directory como URL de redirección, pude empezar a trabajar con llamadas a la API REST para obtener información sobre los calendarios del usuario y sus eventos.\newline
Más adelante vi que nos facilitaría la vida el uso de calendarios compartidos ya que todos serían gestionados desde una cuenta administrador en Outlook y estos pueden ser compartidos con los permisos deseados a los propietarios de las aulas, que serán los encargados de realizar reservas sobre dichos calendarios. De esta manera generamos calendarios globales, que equivalen a aulas, y están disponibles para todos los usuarios con los que se compartan. El problema de ésto lo veremos más adelante, ya que compartir estos calendarios tiene que ser una función a realizar por el administrador manualmente.\newline
Un aspecto relevante del proyecto fue la elección del punto de enlace de Microsoft Graph en vez del extremo de la API de Outlook, ya que ofrece más servicios y características. Elección tomada según la documentación de Microsoft \url{https://docs.microsoft.com/es-ES/outlook/rest/compare-graph}\cite{microsoftGraphCompare}. Primeramente se estuvo trabajando con el punto de conexión de Outlook en la v2.0, ya que parecía más nuevo y posiblemente mejor, pero me fui encontrando errores a la hora de utilizarlo y falta de características, lo que me hizo retroceder y perder tiempo en esto, teniendo que buscar una solución y finalmente encontrando que cambiar los puntos de conexión a Microsoft Graph era la mejor.\newline

Durante el desarrollo fueron surgiendo problemas y diferencias en el uso de las API con el tutorial de Microsoft, ya que está diseñado para Django y no para Flask, aquí tuve que decidir si cambiar el comienzo del proyecto y empezar de nuevo con Django para tener este aspecto solucionado o buscar la alternativa en el lenguaje de Flask, mi idea principal era realizar el proyecto en Flask ya que tenía ganas de aprenderlo pues lo tenía entendido como más rápido para el desarrollo, más fácil de aprender y con alcances del proyecto similares, por lo que la solución tomada fue adaptar el código, esto implicó cambios en el código en el que me fijaba y me obligaba a buscar una adaptación de distintos características que tenían que ser usadas. Por ejemplo, para obtener el JWT (JSON Web Token), no nos sirve la función proporcionada ya que se necesitan diferentes atributos, mediante un identificador de token  proporcionado por una llamada a la API REST en este caso.\newline
Como comentaba anteriormente, un error que no se ha podido solventar es automatizar el compartir calendarios desde la aplicación web, ya que no proporciona esta funcionalidad la API REST de Outlook. Se estuvo buscando en la documentación de la que obtenemos todas las funcionalidades de dicha API REST y lo único que encontramos, fue una funcionalidad en fase Beta aún parte de ella aún y la funcionalidad que nos ofrece es la de gestionar los permisos de un calendario compartido con algún usuario, como vemos en la cita textual de esta página \textit{Antes de que se pueda aplicar el uso compartido del calendario o la delegación, el propietario envía una invitación a una persona con la que compartir o delegado, y dicha persona o delegado acepta la invitación}\cite{compartirCalendario}. Investigando y dando vueltas a este problema, se llegó a la conclusión de que la única forma de compartirlos es manualmente desde la cuenta del administrador.\newline
Otro aspecto a tener en cuenta fue la decisión de subir el proyecto a Azure, ya que se había registrado la aplicación en Azure Active Directory y facilitaba una base de datos y permitía alojar la aplicación web sin coste.\newline
En cuanto al trabajo con Azure, comentar que para realizar la subida tuve que reestructurar el proyecto según lo que Microsoft establece al trabajar en un proyecto que utiliza VSCode y Flask.
\begin{itemize}
    \item \href{https://github.com/microsoft/python-sample-vscode-flask-tutorial}{Estructura VSCode y Flask (Documentación microsoft)}\cite{pythonSample}
    \item \href{https://code.visualstudio.com/docs/python/tutorial-flask}{Estructura VSCode y Flask (Documentación Visual Studio)}\cite{flaskTutorialVSCode}
\end{itemize}
Esto conllevó errores en las partes del código que se importaban al cambiar la carpeta raíz desde VSCode. Siguiendo las mismas especificaciones también generé el fichero \textit{launch} para poder ejecutar desde local mediante el debug.\newline
Una vez que el proyecto estaba estructurado seguí los pasos del tutorial que proporciona Miscrosoft para realizar la subida de la aplicación web a los servidores de Azure. \href{https://docs.microsoft.com/es-es/azure/developer/python/tutorial-deploy-app-service-on-linux-01}{Despliegue en Azure}\cite{deployVSCode}.\newline

Otro aspecto a remarcar es la conexión con la base de datos, al subir todo a los servicios de Azure, la conexión ya no se puede realizar como al principio, que se hacía con la base de datos local de MySQL, ahora se realiza una conexión de pyodbc especificando el servidor, la base de datos, el usuario, contraseña y driver a utilizar desde Azure. Esto no permitirá que se ejecutado en cualquier servidor local como se verá más adelante en la explicación del despliegue. Tuve que modificar las consultas que realizaba ya que no serían desde el mismo cursor ni la misma conexión y el lenguaje de consulta cambia de MySQL al utilizarlo en SQL Database proporcionado por Azure.\newline
Como última anotación, comentar que la parte de la vista para el cliente de la wep app y los estilos que daría a la página los dejo para el final, cuando las funcionalidades estén completadas se mejorará la visualización para el cliente.\newline





\capitulo{6}{Trabajos relacionados}

Esta aplicación web tiene similitudes con todas las aplicaciones de reserva de espacios, pero en lo que distingo mi aplicación principalmente es en el uso de Outlook tanto para la validación del usuario como para la gestión de los calendarios usando Outlook calendar.\newline
En cuanto a funcionalidades, mirando aplicaciones similares de reservas, la principal diferencia a destacar es que ésta está realizada para la universidad o entidades similares, ya que tengo en cuenta distintos edificios en los que se manejan distintas aulas y dichas aulas tienen características y se proporcionan filtros concretos de aulas de universidad.\newline Por ejemplo, aplicaciones similares que he encontrado:
\begin{itemize}
    \item \textit{BookedApp}\newline 
    Es una aplicación que se dedica a la reserva de sitios individuales, no aulas, de biblioteca, mi aplicación web gestiona más de un espacio pero no por sitios individualmente, sino reservas de aulas completas.
    \item  \textit{Booked Scheduler – Reserva de aulas}\href{https://valentingom.wordpress.com/2016/07/10/booked-scheduler-reservas/}{Enlace}\cite{BookedScheduler} \newline
    Es una aplicación más desarrollada que permite reservas de aulas de forma similar a la nuestra, no está sincronizado con Outlook y todo se lleva en base de datos, sin embargo como ventaja suya podemos apreciar los eventos representados en instancias de calendarios.
\end{itemize}
(\textit{BookedApp}) que se dedica a la reserva de sitios de biblioteca, mi aplicación web gestiona más de un espacio pero no por sitios individualmente, sino reservas de aulas completas.\newline


\capitulo{7}{Conclusiones y Líneas de trabajo futuras}

Todo proyecto debe incluir las conclusiones que se derivan de su desarrollo. Éstas pueden ser de diferente índole, dependiendo de la tipología del proyecto, pero normalmente van a estar presentes un conjunto de conclusiones relacionadas con los resultados del proyecto y un conjunto de conclusiones técnicas. 
Además, resulta muy útil realizar un informe crítico indicando cómo se puede mejorar el proyecto, o cómo se puede continuar trabajando en la línea del proyecto realizado. 



\bibliographystyle{plain}
\bibliography{bibliografia}

\end{document}
